\section{Returning the estimate to the chosen family}

% BEGIN P84B1
\paragraph{P84B1. The native good-cut existence theorem.}
For every $D$ and every specified norm choice $a$, assume
$\mathrm{GaloisPreconditions}(D,a)$ and $\mathrm{Tie4}(D)$.
There exists a single cut $c$ such that
\[
  c.\mathrm{Good},\qquad
  \mathrm{NormChoiceFacts}(D_c,a_c),\qquad
  \mathrm{Tie4}(D_c),
  \qquad a_c=\mathrm{nchTr}(c,a).
\]
This is the whole assertion called $\mathrm{CutGivesAlt}$.
Indeed choose $c$ by P84C3. C84F1 and C84F2 transport the
residual and row data and give the cut preconditions and
$\mathrm{CutRec}$ there. C84F3--C84F6 now supply every
norm-choice fact on precisely that cut and transported norm.
T83C3 transports all four ties there. Thus the one choice of
$c$ witnesses the entire conjunction. P84C4 also gives its
commuting condition. No choice in this construction depends on
a subsequent representation, minor, Selmer bound or lift.
% END P84B1

% BEGIN P84B2
\paragraph{P84B2. A sufficient type-prime test for the original family.}
Work under $\mathrm{Hypotheses}_2$ with a given family
$\mathrm{Fm}$ satisfying $\mathrm{TwGood}$. At each original
member $(x,i)$ with $x>\mathrm{Fm}.x_1$, assume the Galois
preconditions for its specified norm choice and the four ties.
Fix one natural number $C_T$ for the family. Suppose that on
one good cut furnished by P84B1 at each such member the Smith
differences of P84D2, divided by $e$, obey the point bound
of P84D3 on every one of its native type-prime tests. Then
$\mathrm{LinkN}(\mathrm{Fm})$ holds. The same conclusion holds
if these differences are bounded directly by the native
$\mathrm{rhs}(C_T,\eta,i,\rho,\pi)$ in every such test.

Choose the specified cuts once for all members. At $x\le
\mathrm{Fm}.x_1$ use the cut retaining all indices and points,
so the cut family is defined everywhere. Its structural maps are
the restrictions of the original maps: its place data, scale,
parameter ledger and reference theta ranks are the original
ones, and its extra norm choice is $a_c$. In particular the
twist test restricts to this family at the original twist
threshold: every retained test for the cut uses the same
character, place and norm comparison in the uncut test.
P84B1, P84C4 and P84D3, or P84D1 in the direct-bound version,
give $\mathrm{FactsFAlt}(\mathrm{gd}\,\mathrm{Fm}\,c,C_T,C_T)$
for all $x>\mathrm{Fm}.x_1$. V84A5 applies at these common
constants and yields $\mathrm{LinkN}$ on the cut family.

For the return to the original family, fix a permitted lift
$j$ and write $d_j=L.dZ0(j)$. Use the original preconditions:
their good-row and theta-scalar clauses and localisation
uniqueness give the disjoint row count and injective theta map
of LN6. The rank clause of their row record identifies that
rank with $d_j$ at every tested twist. The integral operator of
P84C3 on the chosen projector cut, or the operator clause of
goodness for any chosen good cut, bounds the discarded defect
by LN5. Summing over localisations and weighting over places,
LN7 gives
\[
             \eta_{\rm sum}\le\eta'_{\rm sum}+d_j\log N_G .
\]
Here the normalized place weights times the valuation of the
rational integer $N_G$ sum to at most $\log N_G$ by rational
prime factorization; the same multiplicity cap $d_j$ applies
at each place before this summation. The uniform parameter bound
gives $\log N_G\le C_1 x^{C_1}L_{\rm sc}$.
For fixed positive scale exponent, enlarge the common cutoff
so $x^{C_1}\le b$ and $L_{\rm sc}\ge1$, as in LN9.
Since $f,m\ge1$ at a tested lift, the added term is bounded by
$d_j C_1 f m b L_{\rm sc}^2$. Add $C_1$ to the cut-family
constant. The other summand of the polynomial allowance is
nonnegative. The cut preserves the global cyclotomic data,
the $\mathrm{ShaFinOdd}$ guard and the series $\mathrm{vvTerm}$,
so both the summability test and its contribution are literally
the original ones. This proves the claim with constants chosen
before the member and the lift.

On $\mathrm{Des}(\mathcal C,3)$ the same argument works inside
each original $\mathrm{Closed}$ quantifier whenever the two
member hypotheses and the stated integral bound are supplied
there, with $\mathrm{TwGood}$ furnished by the design.
Use the intersection of their original ranges of scale
exponents and the maximum of their transport thresholds if
those bounds are obtained separately, as in A83G4.
This identifies the sufficient arithmetic work on that class.
P84B1 itself holds for every supplied pair $(D,a)$ at its
stated hypotheses.
% END P84B2

% BEGIN P84B3
\paragraph{P84B3. A single integral lifting estimate suffices at the type primes.}
Keep the family, twist, original preconditions, four ties and
published hypotheses of P84B2. In place of its tests on Hecke
indices, suppose one $C_T\in\mathbb N$ supplies the quantitative
potentially crystalline lattice-surjection estimate of T84L1 on
the chosen good cuts. It suffices to supply it on the degree-one
places, common fields and direct or dual lifting tests used
there. Then the original family satisfies $\mathrm{LinkN}$.

Here are the identifications in this implication. P84B1 fixes the
cut and its norm choice; the four transported ties supply the same
depth-zero witness for P84D4 and K84P1. For each good point
P84D4 gives one type locus for every point in its support.
For any native paired test, row admissibility on the cut gives
the base-change, attachment, twist and parameter comparisons of
K84P1. Thus its local compatibility holds at both lifts of each
type place. The rank-one character comparison and annihilating
inertia polynomial for K84P2--K84P4 are exactly the first and
third transported ties.

The period bridges here identify the inertia subgroups: under the
specified isomorphism of decomposition groups inertia is the kernel
of the action on the residue-field closure in both descriptions.
Complete splitting of $g_k$ gives $2f$ degree-one completions.
For a crystalline representation, restricting it further preserves
crystallinity and its labelled jumps. Indeed a crystalline period
basis remains a basis in the period tensor product; invariant
coefficients at the new field are its unramified constants.
For de Rham periods use the new field's constants and tensor the
filtration with that field, giving the jumps at the restricted
embedding. The crystalline tensor and dual comparisons pair the
dual and the original period bases perfectly and the Tate twist
translates the filtration by its one jump. These are the
Fontaine comparisons on the original identified local period
rings. Together with $\mathrm{tiePH}$ from the local bridge,
they are exactly the restriction, inertia and period inputs of
K84P2--K84P4 and T84L5. Consequently there is a direct potential
family on each of our point supports. T84L5 gives the dual
family when required by the sign and obtains both signed jump
bounds from the tested host data.

Now T84L1 uses the assumed quantitative lifting estimate.
T84L2--T84L6, with local finiteness and Bloch--Kato and Tate
duality, give the bounds for the rank block and for every
character. For S84G2 the residual full span at the tested point
is V84A3's residual matrix-span argument: reduce the semisimple
decomposition at that point, cancel its character terms and use
absolute irreducibility and density on the rank block. The
unramified, tame-generator and off-support inputs are those
verified there from the cut datum and C84F4--C84F5. The fourth
coefficient duty gives the original integral idempotent object
as in S84G3. Using global Artin reciprocity and the finiteness
guards of the paired test, S84G3 now bounds precisely the shorter
ideal of P84D5. The finite group contributes at most $2fr^2$.
So the paired condition holds with $C=C_T$ on each of these
cuts. The sum, its return to the uncut family and the uniform
constants are those of P84B2. This implication retains the
integral lifting estimate on the specified surjections; the
existence of potentially crystalline period bases by itself
supplies only the rational comparisons above.
% END P84B3
